\subsection{Classification Results}
\label{sec:wing-results}

Table~\ref{tab:benchmark-main} reports matched AUC, energy independence $I$,
and trainable parameter counts. Within each dataset, all models use the same
metric definitions and evaluation protocol, allowing direct comparison
between specialized baselines and tokenized Transformers. Matched AUC
measures discrimination after balancing the class energy spectra, while $I$
measures within-class score stability with energy.

% Generated by wing_contribution/scripts/build_classification_table.py.
\providecommand{\bestscore}[1]{\ensuremath{\mathbf{#1}}}
\definecolor{datasetbestcolor}{RGB}{0,76,153}
\providecommand{\datasetbest}[1]{\textcolor{datasetbestcolor}{#1}}
\providecommand{\NA}{\textemdash}
\begin{table*}[t]
\centering
\caption{Cross-detector classification results. Higher values are better for energy-matched AUC ($\mathrm{AUC}_{\mathrm{match}}$) and energy independence ($I$). Parameters are rounded counts or ranges in thousands across available dataset-specific implementations; ranges are not uncertainty. Within each model family and dataset, bold marks the highest eligible result where recorded precision resolves the comparison; displayed ties use unrounded values where available. \datasetbest{Blue} marks the highest AUC and $I$ across both model families in each dataset. Highlighting denotes point estimates, not statistical significance. Results are from single training seeds. A dash denotes an unavailable result. $^{\ddagger}$MJD GINE uses additional PSD-label supervision and is excluded from best-result highlighting. Evaluation details are provided in Appendix~\ref{app:metric-details}.}
\label{tab:benchmark-main}
\small
\setlength{\tabcolsep}{2pt}
\renewcommand{\arraystretch}{1}
\begin{tabular*}{\textwidth}{@{\extracolsep{\fill}}lc*{4}{rr}@{}}
\toprule
 & & \multicolumn{2}{c}{NEXT} & \multicolumn{2}{c}{MJD} & \multicolumn{2}{c}{EXO-200} & \multicolumn{2}{c}{SuperNEMO} \\
\cmidrule(lr){3-4}\cmidrule(lr){5-6}\cmidrule(lr){7-8}\cmidrule(l){9-10}
Model & \shortstack{\#Params\\(K)} & $\mathrm{AUC}_{\mathrm{match}}$ & $I$ & $\mathrm{AUC}_{\mathrm{match}}$ & $I$ & $\mathrm{AUC}_{\mathrm{match}}$ & $I$ & $\mathrm{AUC}_{\mathrm{match}}$ & $I$ \\
\midrule
\multicolumn{10}{@{}l}{\textit{Native Representation Approach}} \\
Multi-view CNN & 808 & 0.977 & 0.971 & 0.970 & 0.760 & 0.863 & 0.752 & 0.586 & 0.910 \\
Static GINE & 480 & \datasetbest{\bestscore{0.995}} & 0.970 & 0.969$^{\ddagger}$ & 0.797$^{\ddagger}$ & \bestscore{0.949} & \bestscore{0.801} & \bestscore{0.590} & 0.913 \\
BiGRU & 734 & 0.979 & \bestscore{0.971} & \bestscore{0.975} & \bestscore{0.775} & 0.929 & 0.789 & 0.579 & 0.921 \\
PointMamba-lite & 317 & 0.925 & 0.971 & 0.960 & 0.773 & 0.934 & 0.786 & 0.579 & \bestscore{0.922} \\
\midrule
\multicolumn{10}{@{}l}{\textit{Tokenized Transformers}} \\
Entity + MLP & 111 & 0.975 & \datasetbest{\bestscore{0.975}} & 0.929 & \datasetbest{\bestscore{0.812}} & 0.941 & 0.737 & 0.591 & 0.975 \\
Entity + Fourier & 113--116 & 0.963 & 0.974 & 0.941 & 0.802 & 0.947 & 0.729 & 0.589 & 0.974 \\
Entity + RoPE & 106--113 & 0.497 & 0.512 & 0.793 & 0.749 & 0.946 & 0.715 & 0.551 & 0.865 \\
\addlinespace[2pt]
Region + MLP & 111--121 & \bestscore{0.988} & 0.973 & 0.973 & 0.798 & 0.970 & 0.801 & 0.591 & 0.976 \\
Region + Fourier & 113--125 & 0.979 & 0.973 & \datasetbest{\bestscore{0.981}} & 0.788 & \datasetbest{\bestscore{0.972}} & 0.790 & 0.590 & 0.977 \\
Region + RoPE & 107--116 & 0.512 & 0.483 & 0.916 & 0.768 & 0.971 & 0.789 & \datasetbest{\bestscore{0.602}} & 0.953 \\
\addlinespace[2pt]
Summary + MLP & 111 & 0.987 & 0.971 & 0.965 & 0.799 & 0.947 & 0.776 & 0.590 & 0.975 \\
Summary + Fourier & 113--116 & 0.983 & 0.973 & 0.969 & 0.782 & 0.960 & \datasetbest{\bestscore{0.802}} & 0.589 & \datasetbest{\bestscore{0.977}} \\
Summary + RoPE & 106--114 & 0.602 & 0.532 & 0.878 & 0.777 & 0.955 & 0.773 & 0.582 & 0.956 \\
\bottomrule
\end{tabular*}
\end{table*}

\paragraph{Specialized-baseline comparison.}
Static GINE has the highest matched-AUC point estimate on NEXT, EXO-200,
and SuperNEMO among the shared classic models. It also has the highest $I$
on EXO-200, where the recalculated values are 0.949 and 0.801, respectively.
BiGRU has the highest matched AUC on MJD, while GINE has the highest $I$
with additional PSD-label supervision.
These matched-AUC leaders are unchanged by reevaluation. On SuperNEMO,
PointMamba-lite has the highest $I$ and now marginally exceeds BiGRU in
matched AUC (0.579492745 versus 0.579485174); this small ordering reversal
removes BiGRU from the two-metric Pareto set, leaving GINE and PointMamba-lite.
These single-run point estimates do not establish statistical significance.
Parameter counts and training protocols also
differ, including the additional PSD supervision for MJD GINE.

The finer common energy grid lowers $I$ for the EXO-200 and SuperNEMO
classic models relative to their historical quantile-energy estimates;
both class-level values decline, as documented in
Table~\ref{tab:independence-diagnostics}. This reflects a changed estimator
resolution and pooling rule, not a change in the fitted models.
All reported retained-campaign classic matched AUCs meet the original count and
coverage gates. Coverage, overflow counts, sparse-bin losses, and effective
sample sizes are reported in Appendix~\ref{app:metric-details}.

Among the recorded Transformer results, region tokens provide the strongest
matched-AUC configuration on each detector: region--MLP on NEXT,
region--Fourier on MJD and EXO-200, and region--RoPE on SuperNEMO. Summary
tokens remain close to the leading NEXT representation, suggesting that
spatial aggregation can preserve much of the useful event structure. Entity
tokens are generally weaker on the dense MJD and EXO-200 waveforms, where
selecting isolated measurements can discard contiguous pulse structure
retained by regional patches. Across the completed MLP and Fourier
comparisons, changing tokenization generally produces a larger performance
range than changing positional encoding within one tokenization. This
supports the conclusion that defining what constitutes a token is a central
modeling decision, although the evidence consists of single-seed results.

\paragraph{Effect of positional encoding and energy dependence.}
Positional-encoding preferences remain detector dependent. Coordinate MLPs
perform best across the NEXT tokenizations, whereas Fourier features provide
the strongest MLP/Fourier results on MJD and EXO-200. SuperNEMO shows smaller
differences among the additive encodings, with its reported region--RoPE
variant achieving the highest Transformer matched AUC. NEXT and MJD are more
sensitive to RoPE, but for different reasons: NEXT supplies physical spatial
coordinates, while MJD supplies time together with waveform level and local
change. The comparison therefore varies both coordinate semantics and the
encoding mechanism, so it does not isolate a universal positional-encoding
effect.

Within model subsets evaluated under a common protocol, the strongest
classifier is not always the most energy independent. SuperNEMO further shows
that high independence can coexist with weak class separation. Energy-matched
AUC and $I$ must therefore be interpreted jointly: the former measures
discrimination after energy-spectrum balancing, whereas the latter measures
within-class score stability with energy.
